\chapter{Spot Markets}\label{ch:m3:spot-markets}

The same bitcoin trades on dozens of venues at dozens of prices. Most of the time the prices differ by
a few basis points, the width of the fees; sometimes they do not. From December 2017 to early February
2018, bitcoin bought with won in Seoul cost on average more than 15\% more than bitcoin bought with
dollars in the United States, and for several days 40\% more. Anyone could see the gap, the coins
could cross the world in an hour, and the gap stayed open for weeks. This chapter is about how crypto
spot prices are formed across fragmented venues, the pairs they are quoted in, the arbitrages that
tie them together and what limits those arbitrages, the most important of which is the time and the
permission it takes to move money, and finally about how much of the volume that venues report is real.

\section{Fragmentation and price formation across venues}

A crypto asset has no primary market. There is no listing exchange whose auction sets the reference
price, no consolidated tape and no rule that routes an order to the best price, as in the equity
markets of One Quant Book 1, chapter 9. Each \omterm{def:m3:centralised-exchanges:cex}{centralised exchange} runs its own order book in its own
pairs for its own customers, decentralised venues add their pools (\cref{ch:m3:automated-market-makers}),
and the price of bitcoin is whatever arbitrageurs make these books agree on.

They agree closely where arbitrage capital moves freely. Makarov and Schoar, studying tick data from 34 exchanges in 19
countries, found that price deviations between exchanges within a country typically did
not exceed 1\% on average, while deviations across countries were much larger, and that between
cryptocurrencies on the same venues they were smaller still: when bitcoin in won cost more than 20\% more
than bitcoin in dollars, the price of ether in bitcoin differed by about 3\% on average between the two
countries. Coins moved freely; fiat money did not. The price of a crypto asset is therefore best thought
of as one price per currency area, joined by the cost of moving money between them.

Within a currency area, price discovery happens where the informed trade: in the deepest books, and in the
perpetual futures of \cref{ch:m3:perpetual-futures} that trade alongside spot on the same venues.
Smaller venues follow, their prices kept in line by traders who quote on them against hedges on the large
ones. For a trading firm the practical questions are where the price is made, how quickly others follow,
and what it costs to trade on each venue and to move inventory between them.

\section{Stablecoin quote pairs}

Crypto pairs follow the conventions of foreign exchange (One Quant Book 2, chapter 14): in BTC/USDT
bitcoin is the base currency and USDT the quote currency, and the price is the number of USDT per bitcoin.
Most offshore venues list their deepest books against \omterm{def:m3:blockchains-for-traders:stable}{stablecoins} rather than dollars, because a venue
without banking relationships can hold \omterm{def:m3:blockchains-for-traders:stable}{stablecoins} but not dollar deposits, and because \omterm{def:m3:blockchains-for-traders:stable}{stablecoins} move
between venues on chains at any hour.

\begin{dated}{2026-09}{Quote currencies on one large venue}\label{dat:m3:spot-markets:quotes}
On 24 September 2026 Binance's exchange-information endpoint listed 1{,}372 spot pairs in trading, of
which 495 were quoted in USDT, 264 in USDC, 309 in Turkish lira, 38 in bitcoin and 29 in euros; only a
few were quoted in dollars.
\end{dated}

A price in USDT is a price in dollars only if USDT is worth a dollar. A trader who compares BTC/USDT on one
venue with BTC/USD on another compares two different things, joined by the USDT/USD rate, which is itself
quoted and can move: by a few basis points in normal times, by much more in a \omterm{def:m3:blockchains-for-traders:stable}{depeg}
(\cref{ch:m3:blockchains-for-traders}). The BTC/USDT price and the BTC/USD price give an implied USDT/USD
cross rate, and every arbitrage between the two books is a position in that cross. A firm keeps its books
in one currency and marks its \omterm{def:m3:blockchains-for-traders:stable}{stablecoin} balances at their market price, not at par.

\section{Triangular and cross-venue arbitrage}

Two arbitrages keep prices in line.

\begin{definition}[Triangular arbitrage]\label{def:m3:spot-markets:tri}
\emph{Triangular arbitrage}\index{triangular arbitrage} is a cycle of three trades on one venue that
starts and ends in the same asset through two others, for example USDT into bitcoin, bitcoin into ether and
ether back into USDT, executed when the product of the three exchange rates, after fees, exceeds one.
\end{definition}

\begin{proposition}[The no-arbitrage band of a cross rate]\label{prop:m3:spot-markets:band}
Let a venue quote $A/Q$, $B/Q$ and $B/A$, and charge a taker fee $f$ on each trade. The cycle
$Q \to A \to B \to Q$ pays if and only if
\[
\frac{\text{bid}_{B/Q}}{\text{ask}_{A/Q}\,\text{ask}_{B/A}}\,(1 - f)^3 > 1,
\]
and the reverse cycle $Q \to B \to A \to Q$ pays if and only if
$\text{bid}_{B/A}\,\text{bid}_{A/Q}\,(1 - f)^3 / \text{ask}_{B/Q} > 1$. Neither pays as long as the quoted
$B/A$ price lies within a band around the implied cross $B/Q \div A/Q$ whose half-width is about $3f$
plus the spreads.
\end{proposition}

\begin{proof}
Spend one unit of $Q$: buying $A$ at its ask gives $(1-f)/\text{ask}_{A/Q}$ of $A$, buying $B$ with it at the
ask of $B/A$ gives $(1-f)^2/(\text{ask}_{A/Q}\,\text{ask}_{B/A})$ of $B$, and selling $B$ at its bid returns the
stated product. The reverse cycle is the same computation in the other direction. Taking logarithms, each
condition says that the log of the quoted cross differs from the log of the implied cross by more than
$-3\ln(1-f) \approx 3f$ plus half-spreads.
\end{proof}

At a taker fee of 10 basis points the band is about 30 basis points wide on each side, and a venue's own
market makers keep its cross rates well inside it: triangular opportunities on a single large venue are
rare, small and taken by whoever has the lowest fees and the fastest connection. The tutorial counts them
in synthetic books: 1{,}254 cycles that pay before fees in 500 snapshots of three venues, 26 at a taker fee
of 10 basis points.

\begin{omfigure}
\centering
\begin{tikzpicture}[font=\small,
  c/.style={draw,thick,circle,minimum size=1.3cm,align=center},
  arr/.style={-{Stealth[length=2.2mm]},thick}]
\node[c,draw=omDef,fill=omDef!8] (q) at (0,0) {USDT};
\node[c,draw=omThm,fill=omThm!8] (a) at (3.6,2.2) {BTC};
\node[c,draw=omProp,fill=omProp!8] (b) at (7.2,0) {ETH};
\draw[arr] (q) -- node[above left,font=\footnotesize,align=center] {buy BTC/USDT\\at the ask} (a);
\draw[arr] (a) -- node[above right,font=\footnotesize,align=center] {buy ETH/BTC\\at the ask} (b);
\draw[arr] (b) -- node[below,font=\footnotesize] {sell ETH/USDT at the bid} (q);
\node[font=\footnotesize,align=center] at (3.6,0.75) {three taker fees:\\pays if the product of\\rates exceeds $(1-f)^{-3}$};
\end{tikzpicture}

\omcaption{A triangular cycle on one venue: USDT into bitcoin, bitcoin into ether, ether back into USDT.
The cycle, or its reverse, pays only when the quoted ETH/BTC price leaves the band around the implied cross
set by three taker fees and the spreads. Schematic.}\label{fig:m3:spot-markets:triangle}
\end{omfigure}

\begin{definition}[Cross-venue arbitrage, prepositioned inventory]\label{def:m3:spot-markets:cross}
\emph{Cross-venue arbitrage}\index{cross-venue arbitrage} is buying an asset on one venue and selling it on
another where it is dearer. \emph{Prepositioned inventory}\index{prepositioned inventory} is the stock of
the asset and of the quote currency that an arbitrageur keeps on each venue so that both legs can be
executed at once, the inventory being moved back between venues later, when it is cheapest to do so.
\end{definition}

With inventory on both venues the arbitrageur sells on the dear venue from its stock there and buys on the
cheap one with its cash there, at the same moment; the trade is locked in at the observed prices. The costs
are two taker fees, the eventual cost of rebalancing (a withdrawal fee, a network fee and the time the
inventory spends in transit), and the capital tied up: inventory sitting on venues is exposed to their
failure (\cref{ch:m3:centralised-exchanges}). Without inventory the arbitrageur must buy, transfer, and then
sell, and the price can move while the coins are in flight.

\section{Transfer latency as risk}

\begin{definition}[Transfer latency]\label{def:m3:spot-markets:latency}
\emph{Transfer latency}\index{transfer latency} is the time between the decision to move an asset from one
venue to another and the moment it can be traded at the destination: the source venue's withdrawal
processing, the chain's inclusion and the confirmations the destination requires, and the destination's
crediting.
\end{definition}

The chain's part is set by its block times and \omterm{def:m3:blockchains-for-traders:finality}{finality} (\cref{ch:m3:blockchains-for-traders}); the venues'
parts, which are often longer, by their own processes, reviews and hot-wallet balances. For an arbitrageur
without inventory at the destination, \omterm{def:m3:spot-markets:latency}{transfer latency} turns a locked gap into a bet.

\begin{proposition}[The probability that a transfer arbitrage loses]\label{prop:m3:spot-markets:ploss}
Buy at a price $P$ where the other venue's price is higher by a gap of $g$ (as a fraction of $P$), with
costs $c$, and sell after a transfer of duration $t$ (in years) during which the log price moves by a normal
amount with volatility $\sigma$ and no drift. To first order the trade loses with probability
\[
\Pr(\text{loss}) \approx \Phi\!\left(-\frac{g - c}{\sigma\sqrt{t}}\right),
\]
where $\Phi$ is the standard normal distribution function.
\end{proposition}

\begin{proof}
The net return is $g - c + \varepsilon$ with $\varepsilon \sim \mathcal{N}(0, \sigma^2 t)$ to first order in
the small quantities; it is negative when $\varepsilon < -(g - c)$.
\end{proof}

\begin{omfigure}
\begin{tikzpicture}
\begin{axis}[width=11cm,height=5cm,
  xlabel={transfer time, minutes},ylabel={probability of loss, \%},
  tick label style={font=\small},label style={font=\small},
  xmin=0,xmax=60,ymin=0,ymax=35,grid=major,grid style={gray!25},
  legend columns=3,legend style={font=\footnotesize,at={(0.5,-0.3)},anchor=north,draw=none,fill=none,/tikz/every even column/.append style={column sep=8pt}},
  table/col sep=comma]
\addplot[omThm,very thick] table[x=minutes,y=g50]{figdata/markets-3/16-spot-markets/latency.csv};\addlegendentry{gap 50 bp}
\addplot[omDef,very thick,dashed] table[x=minutes,y=g100]{figdata/markets-3/16-spot-markets/latency.csv};\addlegendentry{gap 100 bp}
\addplot[omExo,very thick,dotted] table[x=minutes,y=g200]{figdata/markets-3/16-spot-markets/latency.csv};\addlegendentry{gap 200 bp}
\end{axis}
\end{tikzpicture}

\omcaption{Probability that a \omterm{def:m3:spot-markets:cross}{cross-venue arbitrage} executed by transferring the asset loses money, as a function
of the transfer time, for three gaps, with 20 basis points of costs and an annual volatility of 60\%
(\cref{prop:m3:spot-markets:ploss}). A 50-basis-point gap loses about one time in three after an hour.
Illustrative parameters. Data: the chapter's tutorial.}\label{fig:m3:spot-markets:latency}
\end{omfigure}

\Cref{fig:m3:spot-markets:latency} is why gaps of tens of basis points survive between venues that are slow to
credit deposits: at 60\% annual volatility the standard deviation of the price over an hour is about 64 basis
points, larger than the gap. The professional answer is \omterm{def:m3:spot-markets:cross}{prepositioned inventory} on every venue that matters, a
rebalancing process that moves it when fees and congestion are low, and a hedge of the inventory itself on a
futures venue. The build of this chapter scans for both kinds of opportunity, charging a trade with transfer the
price risk of its latency and a trade from inventory the cost of rebalancing.

\section{The kimchi premium}

\begin{definition}[Kimchi premium]\label{def:m3:spot-markets:kimchi}
The \emph{kimchi premium}\index{kimchi premium} is the excess of the won price of a crypto asset on South
Korean exchanges over its price elsewhere converted into won at the market exchange rate.
\end{definition}

Makarov and Schoar found that for large parts of 2017 and early 2018 bitcoin on Korean exchanges traded more
than 20\% above the United States, the highest price anywhere; that the daily average price ratio between Korea
and the United States from December 2017 to the beginning of February 2018 was more than 15\% and reached 40\%
for several days; and that the premia opened up in times of large bitcoin appreciation. They estimated the
arbitrage profits available between Korea and the United States over four months at USD~1.275 billion.

The premium lasted because the arbitrage has a fiat leg. The coins can be bought abroad, sent to Korea and sold
for won in an hour; the won must then be converted into dollars and sent back abroad to close the round trip and
start again. That leg was limited: Korean exchanges served residents with local bank accounts, and residents
moving large sums abroad had to justify them. The same authors note that Korean and Japanese exchanges did not
allow short selling, which removed the other way to trade against the premium, selling first and delivering later.

\begin{dated}{2026-09}{Korean outward-remittance documentation}\label{dat:m3:spot-markets:korea}
Makarov and Schoar (2020) describe that in Korea local residents and companies moving more than USD~50{,}000 out
of the country in a single year must submit documents to the authorities proving the reasons for the transfers,
which may not always be approved, and that exchanges operating in several countries let citizens of one country
open accounts only in their local currency. From 30 January 2018 Korea's Financial Services Commission allowed
crypto trading only through real-name bank accounts at the exchange's own bank, closed new accounts to foreigners and
minors, and had banks treat deposits or withdrawals above KRW~10 million a day as suspicious. Rules and thresholds change; a firm checks the current ones with
counsel before trading.
\end{dated}

A capital control does not forbid the arbitrage; it prices it. The weekend problem computes the premium below
which the round trip loses money, given fees, the price risk of the transfer and the cost of getting the won out.

\section{Wash trading and measuring real liquidity}

Volume is how venues are ranked, and a venue's rank brings listings and customers. Many venues have printed
volume that did not happen.

\begin{definition}[Wash trading]\label{def:m3:spot-markets:wash}
\emph{Wash trading}\index{wash trading} is trading in which the same beneficial owner, or parties acting in
concert, are on both sides, so that no risk changes hands; it inflates reported volume and can mislead others
about a market's activity and price.
\end{definition}

In March 2019 the asset manager Bitwise presented to the staff of the SEC an analysis arguing that about 95\% of
the roughly USD~6 billion of daily bitcoin volume reported by data aggregators was fake or non-economic \omterm{def:m3:spot-markets:wash}{wash
trading}, and that the real market was concentrated on ten exchanges whose prices moved tightly together. Cong, Li,
Tang and Yang later applied three statistical tests to 29 exchanges. Trades on regulated exchanges showed the
patterns of trading everywhere: first significant digits of trade sizes following Benford's law, clustering of
sizes at round numbers, and fat-tailed size distributions. Unregulated exchanges often failed them, and the authors
estimated that wash trades averaged over 70\% of the reported volume of the unregulated exchanges.

The two tests the tutorial implements are simple. Benford's law says that in many naturally occurring data the
first significant digit is $d$ with probability $\log_{10}(1 + 1/d)$: a 1 about 30\% of the time, a 9 under 5\%.
Machine-generated sizes drawn uniformly from a range do not follow it. And people trade round sizes, 0.01 or 0.1
bitcoin, far more often than 0.0137; a bot printing volume does not bother. Cong and co-authors used the regulated
exchanges' ratio of unrounded to rounded trades as a benchmark: rounded trades on another exchange, scaled by that
ratio, estimate its authentic trades, and the excess of unrounded trades is wash.

\begin{omfigure}
\begin{tikzpicture}
\begin{axis}[width=11cm,height=5cm,ybar,bar width=6pt,
  xlabel={first significant digit of the trade size},ylabel={share of trades, \%},
  tick label style={font=\small},label style={font=\small},
  xmin=0.4,xmax=9.6,xtick={1,2,3,4,5,6,7,8,9},ymin=0,ymax=35,grid=major,grid style={gray!25},
  legend columns=3,legend style={font=\footnotesize,at={(0.5,-0.3)},anchor=north,draw=none,fill=none,/tikz/every even column/.append style={column sep=8pt}},
  table/col sep=comma]
\addplot[fill=gray!40,draw=gray,area legend] table[x=digit,y=benford]{figdata/markets-3/16-spot-markets/digits.csv};\addlegendentry{Benford's law}
\addplot[fill=omDef!60,draw=omDef,area legend] table[x=digit,y=genuine]{figdata/markets-3/16-spot-markets/digits.csv};\addlegendentry{genuine tape}
\addplot[fill=omThm!60,draw=omThm,area legend] table[x=digit,y=mixed]{figdata/markets-3/16-spot-markets/digits.csv};\addlegendentry{70\% wash trades}
\end{axis}
\end{tikzpicture}

\omcaption{First significant digits of 20{,}000 trade sizes: a synthetic genuine tape follows Benford's law
(chi-square 1.4 against a 5\% critical value of 15.5); a tape in which 70\% of trades are machine-generated with
sizes uniform between 0.0010 and 0.0500 bitcoin does not. Synthetic tapes. Data: the chapter's tutorial.}\label{fig:m3:spot-markets:benford}
\end{omfigure}

\begin{omfigure}
\begin{tikzpicture}
\begin{axis}[width=11cm,height=5cm,
  xlabel={trade size, BTC (bins of 0.001)},ylabel={share of trades, \%},
  tick label style={font=\small},label style={font=\small},
  xmin=0,xmax=0.05,ymin=0,ymax=8,grid=major,grid style={gray!25},
  xticklabel style={/pgf/number format/fixed,/pgf/number format/precision=2},
  scaled x ticks=false,
  legend columns=2,legend style={font=\footnotesize,at={(0.5,-0.3)},anchor=north,draw=none,fill=none,/tikz/every even column/.append style={column sep=8pt}},
  table/col sep=comma]
\addplot[omDef,very thick,const plot] table[x=size,y=genuine]{figdata/markets-3/16-spot-markets/sizes.csv};\addlegendentry{genuine tape}
\addplot[omThm,very thick,dashed,const plot] table[x=size,y=mixed]{figdata/markets-3/16-spot-markets/sizes.csv};\addlegendentry{70\% wash trades}
\end{axis}
\end{tikzpicture}

\omcaption{Trade sizes up to 0.05 bitcoin: genuine trading clusters at round sizes (the bins starting at 0.01,
0.02, 0.03 and 0.04 bitcoin); machine-generated volume dilutes the clusters. Synthetic tapes. Data: the chapter's
tutorial.}\label{fig:m3:spot-markets:sizes}
\end{omfigure}

Reported volume is therefore a poor measure of where to trade. What matters to a firm that must execute is what it
can trade without moving the price.

\begin{definition}[Quoted depth]\label{def:m3:spot-markets:depth}
\emph{Quoted depth}\index{quoted depth} is the quantity resting in an order book within a stated distance of the
mid price, for example within 10 basis points on each side, summed over price levels and expressed in the quote
currency.
\end{definition}

Depth can be faked too, by orders that are cancelled before anyone reaches them, so a firm measures it over time,
checks it against its own fills, and compares venues by what it costs to execute its own sizes: the realised
slippage of its trades, the share of its orders filled at the quoted price, and how quickly the book refills after
a trade. These measurements, not reported volume, decide where the firm sends its orders and where it keeps
inventory.

\section{Tutorial: arbitrage scans and wash-trading tests}\label{tut:m3:spot-markets:scan}

\begin{tutorial}
\textbf{Goal.} Scan synthetic books on three venues for triangular and cross-venue opportunities after fees, price
the risk of \omterm{def:m3:spot-markets:latency}{transfer latency}, and apply two wash-trading tests to synthetic tapes. \textbf{End state:}
\cref{fig:m3:spot-markets:latency,fig:m3:spot-markets:benford,fig:m3:spot-markets:sizes} and the counts in the
text.
\begin{tutsteps}
\item \textbf{The triangle.} Both directions of a cycle, with the executable size at the top of the books.
\omcode{code/firm/triarb/firm_triarb.py}{30}{43}{Both directions of a triangular cycle on one venue.}\label{lst:m3:spot-markets:tri}
\item \textbf{The scan.} \texttt{count\_opportunities(fee)} runs the scanner on 500 snapshots at taker fees of 0,
2, 5 and 10 basis points.
\item \textbf{The tests.} First digits against Benford's law, clustering at multiples of 0.01 bitcoin, and the
rounding-based estimate of the wash share.
\omcode{code/markets-3/16-spot-markets/python/m3_spot.py}{93}{113}{The clustering ratio and the rounding-based estimate of the share of wash trades.}\label{lst:m3:spot-markets:wash}
\item \textbf{Run} \texttt{fig\_spot.py} for the chart data.
\end{tutsteps}
\textbf{What to change next.} Let the wash trader round its sizes and see which test still catches it; add a fat-tail
test on the largest trades; make the synthetic venues' prices follow one leader with a lag and measure which venue
leads.
\end{tutorial}

\section{Build: the arbitrage scanner}\label{bld:m3:spot-markets:triarb}

\begin{build}
\bfield{Purpose} Find, across the venues the miniature firm is connected to, the triangular cycles and cross-venue
trades that pay after every cost, and rank them by expected profit; feed the execution layer and the inventory
planner.
\bfield{Interface} \texttt{Quote(bid, ask, bid\_size, ask\_size)}; \texttt{triangular(book, a, b, q, fee)} returns
both directions; \texttt{cross(buy, sell, fee\_buy, fee\_sell, route, rebalance\_bp, transfer\_minutes, vol\_annual,
z)}; \texttt{scan(books, fees, cycles, rebalance\_bp, min\_edge\_bp)} returns opportunities, largest profit first.
\bfield{Rules} Taker fees per venue (from each venue's tier); sizes limited by the top of every book on the route;
trades from inventory charged the rebalancing cost, trades with transfer charged $z$ standard deviations of the
move over the transfer time (from \texttt{firm.gasfee}'s confirmation times plus the venues' processing).
\bfield{Acceptance tests} \texttt{code/firm/triarb/tests/}: both triangular directions against hand computation; a
fee that removes an opportunity; the transfer-risk charge; ordering and filtering of the scan.
\bfield{Stretch} Depth beyond the top of the book; \omterm{def:m3:blockchains-for-traders:stable}{stablecoin} cross rates marked at market rather than par; inventory
limits per venue and asset; the scanner run on live books with its hit rate measured against fills.
\end{build}

\begin{omsources}
\item I.~Makarov and A.~Schoar, ``Trading and arbitrage in cryptocurrency markets'', \emph{Journal of Financial
Economics} 135(2), 2020, 293--319 (accepted version, LSE Research Online).
\item L.~W.~Cong, X.~Li, K.~Tang and Y.~Yang, ``Crypto Wash Trading'', NBER Working Paper 30783, 2022; published in
\emph{Management Science}, 2023.
\item SEC, memorandum of a meeting with Bitwise Asset Management, NYSE Arca and Vedder Price, File No.\
SR-NYSEArca-2019-01, 20 March 2019, with the Bitwise presentation.
\item Binance, exchange-information endpoint of the spot API, 24 September 2026.
\item Financial Services Commission (Korea), press release on real-name cryptocurrency trading, 23 January 2018.
\end{omsources}

\section{Exercises}

\begin{exercise}[$\star$]\label{exo:m3:spot-markets:1}
BTC/USDT is 60{,}000 and ETH/USDT 3{,}000. What is the implied ETH/BTC cross rate? At a taker fee of 10 basis
points, roughly how far can the quoted ETH/BTC price be from it before a triangle pays?
\end{exercise}

\begin{exercise}[$\star$]\label{exo:m3:spot-markets:2}
BTC/USDT is 60{,}030 on one venue and BTC/USD 60{,}000 on another. What USDT/USD rate makes the two prices
consistent?
\end{exercise}

\begin{exercise}[$\star$]\label{exo:m3:spot-markets:3}
Why can a bitcoin price gap between two countries be much larger than the gap in the ether--bitcoin price between
the same countries?
\end{exercise}

\begin{exercise}[$\star\star$]\label{exo:m3:spot-markets:4}
A gap of 100 basis points, costs of 20, volatility 60\% a year. What is the probability of loss if the transfer
takes 30 minutes? 60 minutes?
\end{exercise}

\begin{exercise}[$\star\star$]\label{exo:m3:spot-markets:5}
On a benchmark exchange 40\% of trades of at least 0.01 bitcoin are at multiples of 0.01. On another, 12\% are.
Estimate its share of wash trades among such trades.
\end{exercise}

\begin{exercise}[$\star\star$]\label{exo:m3:spot-markets:6}
Explain why \omterm{def:m3:spot-markets:cross}{prepositioned inventory} removes transfer risk from an arbitrage but not all of its risks.
\end{exercise}

\begin{exercise}[$\star\star\star$]\label{exo:m3:spot-markets:7}
\emph{Coding.} With \texttt{count\_opportunities}, find how many triangular and cross-venue opportunities pay in 500
snapshots at taker fees of 0, 2, 5 and 10 basis points. Why do cross-venue opportunities outnumber triangles?
\end{exercise}

\begin{exercise}[$\star\star\star$]\label{exo:m3:spot-markets:8}
\emph{Find the flaw.} ``This exchange reports twice the volume of that one, so our orders will cost less there.''
\end{exercise}

\section{Problem: The Kimchi Premium}

\begin{problem}[{Weekend problem --- pricing a capital control}]\label{pb:m3:spot-markets:1}
A resident of Korea with accounts abroad considers the round trip: buy bitcoin with dollars abroad (taker fee
0.10\%), withdraw it (0.05\% of value), transfer it in an hour, sell it for won in Korea (taker fee 0.05\%), convert
the won to dollars (spread 0.5\%) and remit them abroad through a channel that costs 2\% all-in. Volatility is 70\%
a year; the trader charges the transfer 1.65 standard deviations of the price move. All costs are illustrative.

\medskip\noindent\textbf{Part I --- The costs.}
\begin{enumerate}
\item What does a dollar of bitcoin cost, delivered to the withdrawal, abroad?
\item What is the price-risk charge for the transfer?
\item What fraction of the won proceeds comes back as dollars abroad?
\item What premium makes the round trip break even?
\item At the 15\% average premium of December 2017 to early February 2018, what is the expected profit per dollar,
without the risk charge?
\end{enumerate}

\medskip\noindent\textbf{Part II --- The constraint.}
\begin{enumerate}[resume]
\item What would the break-even premium be if the won could be remitted at no cost beyond the spread?
\item And if the transfer took ten minutes instead of an hour?
\item On USD~50{,}000, what is the expected profit of one round trip at 15\%?
\item Why can a foreign firm not simply sell bitcoin on a Korean exchange?
\item Why could no one sell short in Korea against the premium?
\end{enumerate}

\medskip\noindent\textbf{Part III --- The evidence.}
\begin{enumerate}[resume]
\item What happened to the premium when bitcoin rose fast, and why?
\item Why was the ether--bitcoin gap between Korea and the United States only about 3\%?
\item What do the estimated USD~1.275 billion of arbitrage profits between Korea and the United States tell us?
\item Would \omterm{def:m3:spot-markets:cross}{prepositioned inventory} in Korea remove the constraint?
\item Would moving the value out as a \omterm{def:m3:blockchains-for-traders:stable}{stablecoin} instead of won escape the constraint?
\end{enumerate}

\medskip\noindent\textbf{Part IV --- Judgement.}
\begin{enumerate}[resume]
\item Is the \omterm{def:m3:spot-markets:kimchi}{kimchi premium} an inefficiency?
\item Who earned it, and what did they risk?
\item What does the premium tell a firm about where prices are made?
\item State the \emph{named result}: the break-even premium of the round trip with and without the remittance
cost.
\item In one sentence: what limits arbitrage between crypto venues?
\end{enumerate}
\end{problem}

\section{Interview questions}

\begin{interviewq}[$\star$ \iqroles{trader}]\label{iq:m3:spot-markets:1}
What is \omterm{def:m3:spot-markets:tri}{triangular arbitrage}, and why is it rare on a large venue?
\end{interviewq}

\begin{interviewq}[$\star$ \iqroles{trader, risk}]\label{iq:m3:spot-markets:2}
Why is a BTC/USDT price not a dollar price, and what does that mean for your books?
\end{interviewq}

\begin{interviewq}[$\star\star$ \iqroles{researcher}]\label{iq:m3:spot-markets:3}
How would you test whether an exchange's reported volume is real?
\end{interviewq}

\begin{interviewq}[$\star\star$ \iqroles{trader}]\label{iq:m3:spot-markets:4}
Two venues differ by 50 basis points for an hour. Why might nobody be arbitraging it?
\end{interviewq}

\begin{interviewq}[$\star\star$ \iqroles{risk}]\label{iq:m3:spot-markets:5}
What risks does an arbitrage desk's \omterm{def:m3:spot-markets:cross}{prepositioned inventory} carry, and how would you limit them?
\end{interviewq}

\begin{interviewq}[$\star\star\star$ \iqroles{developer, researcher}]\label{iq:m3:spot-markets:6}
Design a system that decides where to keep inventory across ten venues for a \omterm{def:m3:spot-markets:cross}{cross-venue arbitrage} strategy.
\end{interviewq}
